\chapter*{ABSTRACT}
\addcontentsline{toc}{chapter}{ABSTRACT}
\justifying
Intelligent Transportation Systems increasingly use federated learning (FL) to let vehicles and infrastructure jointly train traffic-prediction models without sharing raw data. However, data locality alone does not guarantee reliability: faulty or malicious agents may submit false traffic reports or poisoned model updates. The problem is especially important in traffic settings because agents genuinely observe different parts of the road. Disagreement between honest agents is therefore expected and must not be mistaken for misbehaviour. Existing FL security and vehicular-trust techniques address parts of this problem, but a coordinated approach is needed to distinguish legitimate disagreement from attacks and to connect trust with both learning and traffic-state consensus.

The overall project addresses this problem through TRAF-MA, a trust-aware multi-agent architecture for a signalised intersection. Its intended design combines heterogeneous vehicle and traffic-signal agents, lightweight local models for congestion classification and queue-length regression, a shared feature representation, and an edge coordinator. The planned trust engine uses contextual, physical and prediction consistency, model-update quality, network reliability and temporal behaviour. Relevance weighting is intended to account for spatial, temporal and coverage differences. Trust is planned to inform both federated aggregation and traffic-state consensus, and to support anomaly handling, quarantine and recovery under false-observation and model-poisoning attacks. The project also aims to evaluate this system in SUMO and integrate it with a real-time monitoring dashboard.

This report is an interim implementation account toward those continuing project objectives. The work completed so far includes the SUMO/TraCI intersection and bounded agents, future queue and congestion labels, compact dataset preparation, a two-head GRU, baseline and central training, and hierarchical FedAvg with an RSU self-update ablation. In the saved 100-vehicle experiment, the central GRU achieved MAE 18.51 m, RMSE 28.28 m, macro-F1 0.764 and balanced accuracy 0.746 on 300 overlapping held-out windows. Hierarchical FL achieved macro-F1 0.342 with RSU learning enabled and 0.298 with it disabled. These are single-seed results from the current simulation; they do not establish the effectiveness of the planned trust, attack, recovery or dashboard components. The report therefore preserves the overall project problem and objectives while presenting methodology and results for the implementation completed to date.

\textbf{Keywords:} Federated Learning, Intelligent Transportation Systems, Trust Management, Multi-Agent Systems, Traffic Prediction, Traffic-State Consensus, SUMO Simulation.
